\documentclass[main.tex]{subfiles}
\begin{document}
Quyển sách thứ hai trong sê-ri \textit{hologra} (holo no Graffiti) tiếp nối hành trình đầy hài hước của Hikawa Kuon tại văn phòng \hll. Sau năm tuần đầu tiên đầy thử thách và hỗn loạn, mùa hè đầu tiên của cậu tại văn phòng đã chính thức bắt đầu, mang theo những sắc màu rực rỡ và những sự kiện mới không thể ngờ tới. Đây là khoảng thời gian mà cả văn phòng không chỉ lộn xộn về công việc, mà còn trở thành một chiến trường của những trò đùa nghịch, những kế hoạch kỳ quặc, và vô số câu chuyện dở khóc dở cười của các thành viên trong công ty.

Mùa hè cũng đánh dấu sự xuất hiện của những gương mặt mới, mỗi người đều góp phần tạo nên những chuỗi hỗn loạn mới mẻ và thú vị. Từ những buổi họp chẳng đi đến đâu vì bị chen ngang bởi ý tưởng đi chơi biển, cho đến những trận chiến bất tận với chiếc điều hòa "không biết nghe lời", mỗi ngày trong văn phòng dường như là một cuộc phiêu lưu nhỏ.

Không khí mùa hè như khiến mọi người trong công ty càng thêm bùng nổ, biến những tình huống vốn đã hài hước trở nên sôi động và thậm chí là "náo loạn" hơn bao giờ hết. Hikawa Kuon không chỉ phải xoay sở với công việc của mình mà còn phải đối mặt với sự "hợp tác" kỳ lạ từ đồng nghiệp, những người luôn tìm cách kéo cậu vào các rắc rối mới.

Trong quyển số Một này, chúng ta hãy cùng tiếp tục theo chân Kuon và các cô gái cng hòa mình vào "Những Ngày Hè Náo Loạn", nơi mùa hè rực rỡ không chỉ mang theo ánh nắng, mà còn làm bừng lên những câu chuyện vui nhộn không hồi kết với mười ba câu chuyện nhỏ khác nhau!

% Tiêu đề quyển: Những Ngày Hè Náo Loạn

Đúng như tiêu đề, quyển số Một sẽ gồm các thành viên mới sau:
\begin{itemize}
  \item Ookami Mio từ GAMERS.
  \item Aki Rosenthral và Yozora Mel từ Gen 1.
\end{itemize}

Lưu ý về cuốn sách này: Đại da số các tập trong cuốn này (trừ {\flaregothic ÉPISODE 8}) đều đã bị thất lạc, tức không thể tìm thấy trên kênh chính thức, mà chỉ còn sót lại trên kênh \textit{bilibili} cùa \hll. Cuốn này có thể được coi là một cuốn truyện ``thất lạc''.

\underline{Quyển này sẽ bao gồm:}
\begin{itemize}
  \item 12 tập bình thường (đánh số từ 5 đến 16);
  \item 01 tập đặc biệt (đánh số La Mã I);
  \item 01 tập tự làm (tập số Một).
\end{itemize}

\end{document}